\section{Conclusion and Future Work}
\label{sec:conclusion}

This work set out to establish whether a dynamic, template-based benchmark for LLM unit-test generation resists contamination-induced score inflation more effectively than a decontaminated static alternative, without losing the ability to measure genuine test-generation capability. The clearest evidence obtained remains the causal result reported in Section~\ref{sec:results-h2}. The observational comparisons are more equivocal -- direction-consistent but correction-dependent -- an expected property of comparing nine models of unknown training history rather than a controlled intervention. The dynamic benchmark also tracks an external code-generation proxy at least as well as the static benchmark does, addressing the risk that contamination resistance came at the cost of construct validity. Taken together, the six pre-registered hypotheses support dynamic, template-based generation as a practical, contamination-resistant evaluation strategy for LLM test-generation research, while identifying specifically where its present statistical power -- principally in ranking stability across independent re-instantiations -- falls short of a pre-registered bar, rather than leaving that shortfall undiagnosed.

Several directions were scoped out of this study by an explicit, pre-registered cut-priority order rather than omitted incidentally, and remain open: a sensitivity sweep at five attempts per problem and a one-shot prompting condition, both cut first under budget pressure with no experimental work begun; and a fifth independent benchmark instantiation for the two models that did not reach it, whose cost of completion was quantified directly and declined on diagnostic grounds (Section~\ref{sec:results-h5}) rather than left unexamined. Semantic-preserving-change mutation operators, complementary to the semantic-altering-change operators used throughout this study, would extend the mutation-testing methodology to a robustness dimension not evaluated here. A supplementary side study begun during this project, characterising the specific patterns of mutation-explanation failure across models, was not carried through to a reportable result and remains future work in its own right. Finally, public release of the template corpus and the instantiation engine -- balancing the benchmark's contamination resistance, which depends on the template grammar remaining unpublished, against the reproducibility value of releasing the engine that consumes it -- is under consideration pending confirmation from the supervising academic on the appropriate release scope.
